\documentclass[11pt]{article}
\usepackage[T1]{fontenc}
\usepackage{amsmath,amssymb,amsthm,mathtools}
\usepackage[a4paper,margin=1in]{geometry}
\usepackage[hidelinks]{hyperref}
\usepackage{microtype}

\newtheorem{theorem}{Theorem}
\newtheorem{corollary}[theorem]{Corollary}
\newtheorem{lemma}[theorem]{Lemma}
\newtheorem{remark}[theorem]{Remark}
\newcommand{\F}{\mathbb F}
\newcommand{\E}{\mathbb E}
\newcommand{\Prb}{\mathop{\rm Pr}}
\newcommand{\1}{\mathbf 1}

\title{Graphlike Overlap--Peierls Reduction for Frontier}
\author{Research note}
\date{28 July 2026}

\begin{document}
\maketitle

\begin{abstract}
We complete the graphlike family step in the comparison-spectrum analysis of
Frontier pruning.  An explicit subcoupling lifts quotient boundary overlaps
back to prefix assignments and bounds each comparison by one product-law
overlap on the union of its visible connected components.  A two-bit
counterexample shows that exact overlap weights cannot in general be
multiplied componentwise.  The Hellinger majorant does factorize.  Combined
with Finner's inequality, detector-column weight at most two, and score
exponents at most one, it proves the required pointwise polymer domination
and yields fractional top-\(K\), score-gap, connected-count, and integrated
ordering-lifetime bounds.
\end{abstract}

The extension to arbitrary bounded detector-column weight, together with an
exact retained-matrix profile, is given in
\texttt{\detokenize{BOUNDED_HYPERGRAPH_OVERLAP.tex}}.

\section{Prefix quotient and open components}

Fix a cut \(t\), let \(I_t\) be the processed variables, and write
\[
 P_t(u)=\Prb[X_{I_t}=u].
\]
Define the linear quotient map and its pushforward mass by
\[
\begin{split}
 F_t(u)
 &=
 \left(
 H_{C_t,I_t}u,\,
 H_{\Gamma_t,I_t}u,\,
 L_{I_t}u
 \right)
 =(c(u),b(u)),\\
 A_t(y)&=\sum_{u:F_t(u)=y}P_t(u).
\end{split}
\tag{1}
\]
For a nonzero relative boundary shift
\[
 \xi\in G_t=D_t(\ker H_{C_t,I_t}),\qquad
 D_tv=(H_{\Gamma_t,I_t}v,L_{I_t}v),
\]
let \(\delta_t\xi\) be its active-detector component.  At fixed future
active parity \(q\), set
\[
 r_t(q,\xi)=\frac{g_t(q+\delta_t\xi)}{g_t(q)}
\tag{2}
\]
and
\[
 O_{t,q}(\xi)
 =
 \sum_y\min\{A_t(y),r_t(q,\xi)A_t(y+\xi)\}.
\tag{3}
\]
The comparison spectrum of the companion note is
\[
 \omega_t(\xi)=\E_{Q_t}O_{t,Q_t}(\xi).
\tag{4}
\]

Let \(\mathcal G_t\) be the graph on \(I_t\) in which two variables are
adjacent when they both occur in a completed detector row.  A nonempty
connected set \(\gamma\) is an open-prefix polymer when
\[
 H_{C_t,I_t}\1_\gamma=0.
\tag{5}
\]
It is visible when \(D_t\1_\gamma\ne0\).  Polymers are compatible when no
graph edge joins them.  For a compatible family \(\mathcal F\), define
\[
 S_{\mathcal F}=\bigcup_{\gamma\in\mathcal F}\gamma,\qquad
 D_t\mathcal F
 =
 \bigoplus_{\gamma\in\mathcal F}D_t\1_\gamma.
\tag{6}
\]

\begin{lemma}[componentwise completed parity]
If \(v\in\ker H_{C_t,I_t}\), every connected component \(\gamma\) of
\(\operatorname{supp}v\) in \(\mathcal G_t\) satisfies
\(H_{C_t,I_t}\1_\gamma=0\).
\end{lemma}

\begin{proof}
A completed row cannot meet two different support components, because any
two support variables in that row are adjacent.  The row has even parity on
\(v\), hence even parity on the unique component it meets.
\end{proof}

\section{Exact coupling lift}

For \(S\subseteq I_t\) and \(r>0\), define
\[
 \kappa_{t,r}(S)
 =
 \sum_{x\in\F_2^S}
 \min\{P_{t,S}(x),rP_{t,S}(x+\1_S)\},
\tag{7}
\]
where \(P_{t,S}\) is the product marginal.  Put
\[
 \overline\kappa_t(\mathcal F)
 =
 \E_{Q_t}
 \kappa_{t,r_t(Q_t,D_t\mathcal F)}(S_{\mathcal F}).
\tag{8}
\]

\begin{theorem}[quotient coupling lift]
For every nonzero \(\xi\in G_t\),
\[
\boxed{
 \omega_t(\xi)
 \le
 \sum_{\substack{
 \mathcal F\ {\rm compatible\ visible}\\
 D_t\mathcal F=\xi}}
 \overline\kappa_t(\mathcal F).}
\tag{9}
\]
\end{theorem}

\begin{proof}
Fix \(q\), abbreviate \(r=r_t(q,\xi)\), and let
\[
 \nu_y=\min\{A_t(y),rA_t(y+\xi)\}.
\]
For nonzero fibers, define a flow between prefix assignments:
\[
 \lambda_y(u,u')
 =
 \nu_y
 \frac{P_t(u)}{A_t(y)}
 \frac{P_t(u')}{A_t(y+\xi)},
 \qquad F_t(u)=y,\quad F_t(u')=y+\xi.
\tag{10}
\]
Its total mass is \(O_{t,q}(\xi)\), its first marginal is at most \(P_t\),
and its second marginal is at most \(rP_t\).

Every pair carrying flow has
\[
 v=u+u'\in\ker H_{C_t,I_t},\qquad D_tv=\xi.
\tag{11}
\]
Partition the flow by the compatible family \(\mathcal F\) of visible
connected components of \(\operatorname{supp}v\), ignoring invisible
components.  For \(S=S_{\mathcal F}\), every pair in that class satisfies
\[
 u'_S=u_S+\1_S.
\]
After projection to \(S\), the flow from \(x\) is at most both
\(P_{t,S}(x)\) and \(rP_{t,S}(x+\1_S)\).  The class therefore has mass at
most \(\kappa_{t,r}(S)\).  Sum over families and average over \(Q_t\).
\end{proof}

\begin{remark}
The lift handles arbitrary quotient merging and arbitrary invisible kernel
components.  It does not choose one representative for a boundary shift.
\end{remark}

\section{Exact overlap does not factorize}

Consider two independent Bernoulli-\(p\) processed variables, no completed
constraints, the identity boundary map, and \(g_t\equiv1\).  The shift that
flips both variables consists of two compatible visible singleton polymers.
For \(0<p<1/2\),
\[
 \Omega_2(p)=2p,
 \qquad
 \Omega_1(p)^2=(2p)^2=4p^2<2p.
\tag{12}
\]
Thus exact componentwise overlap activities do not multiply, even before
quotient merging.  The correct exact activity in Theorem 2 is one overlap on
the family union.  In the homogeneous unscored model,
\[
\boxed{
 \overline\kappa_t(\mathcal F)
 =
 \Omega_{|S_{\mathcal F}|}(p),}
\tag{13}
\]
where, for \(J\sim{\rm Bin}(m,p)\),
\[
 \Omega_m(p)
 =
 2\Prb[J>m/2]+\Prb[J=m/2].
\tag{14}
\]

Let \(M_t(m)\) count compatible visible families with nonzero total boundary
shift and union size \(m\).  For \(0<\rho<1\), subadditivity gives the exact
finite-size family bound
\[
\boxed{
 \sum_{\xi\ne0}\omega_t(\xi)^\rho
 \le
 \sum_{m\ge1}M_t(m)\Omega_m(p)^\rho.}
\tag{15}
\]
The sharp cap theorem in the companion note therefore gives
\[
\boxed{
 D_t^{\rm cap}
 \le
 K^{-1/\rho}
 \left[
 \sum_{m\ge1}M_t(m)\Omega_m(p)^\rho
 \right]^{1/\rho}.}
\tag{16}
\]

\section{Hellinger factorization and graphlike theorem}

For processed variable \(j\), define
\[
 \beta_j=2\sqrt{p_j(1-p_j)}.
\tag{17}
\]

\begin{lemma}[tilted Hellinger majorant]
For \(S\subseteq I_t\) and \(r>0\),
\[
\boxed{
 \kappa_{t,r}(S)
 \le
 \sqrt r\prod_{j\in S}\beta_j.}
\tag{18}
\]
\end{lemma}

\begin{proof}
Apply \(\min\{a,rb\}\le\sqrt r\sqrt{ab}\) in (7).  The remaining
Hellinger coefficient factorizes over the product marginal.
\end{proof}

Define the score moment
\[
 R_{t,1/2}(\delta)
 =
 \E_{Q_t}
 \left[
 \frac{g_t(Q_t+\delta)}{g_t(Q_t)}
 \right]^{1/2}.
\tag{19}
\]
Then
\[
 \overline\kappa_t(\mathcal F)
 \le
 R_{t,1/2}(\delta_tD_t\mathcal F)
 \prod_{j\in S_{\mathcal F}}\beta_j.
\tag{20}
\]
For the row-factorized score
\[
 g_t(q)
 =
 \prod_{i\in\Gamma_t}
 \rho_{i,t}(q_i)^{\alpha_{i,t}},
 \qquad
 \rho_{i,t}(b)=\Prb[(Q_t)_i=b],
\tag{21}
\]
Finner's inequality gives
\[
 \frac12
 \max_{j\in J_t}
 \sum_{\substack{i:\delta_i=1\\H_{ij}=1}}\alpha_{i,t}
 \le1
 \quad\Longrightarrow\quad
 R_{t,1/2}(\delta)\le1.
\tag{22}
\]

\begin{theorem}[graphlike pointwise domination]
Suppose every detector column has weight at most two, the score has form
(21), and \(\alpha_{i,t}\le1\) for every active row.  Define
\[
 w_t(\gamma)=\prod_{j\in\gamma}\beta_j.
\tag{23}
\]
Then, for every nonzero \(\xi\in G_t\),
\[
\boxed{
 \omega_t(\xi)
 \le
 \sum_{\substack{
 \mathcal F\ {\rm compatible\ visible}\\
 D_t\mathcal F=\xi}}
 \prod_{\gamma\in\mathcal F}w_t(\gamma).}
\tag{24}
\]
\end{theorem}

\begin{proof}
A future detector column meets at most two active rows.  The left side of
(22) is therefore at most \(\frac12(2)=1\).  Hence the score moment in
(20) is at most one.  Compatible components are disjoint, so
\[
 \prod_{j\in S_{\mathcal F}}\beta_j
 =
 \prod_{\gamma\in\mathcal F}w_t(\gamma).
\]
Apply Theorem 2.
\end{proof}

\section{Fractional Peierls and score-gap bounds}

Let \(\mathcal P_t^{\rm vis}\) be the visible polymers and define
\[
 \Xi_{t,\rho}
 =
 \sum_{\gamma\in\mathcal P_t^{\rm vis}}w_t(\gamma)^\rho.
\tag{25}
\]
For \(0<\rho<1\), Theorem 4 and subadditivity imply
\[
\boxed{
 \sum_{\xi\ne0}\omega_t(\xi)^\rho
 \le
 \prod_{\gamma\in\mathcal P_t^{\rm vis}}
 (1+w_t(\gamma)^\rho)-1
 \le e^{\Xi_{t,\rho}}-1.}
\tag{26}
\]
Consequently,
\[
\boxed{
 D^{\rm cap}
 \le
 K^{-1/\rho}\sum_t
 \left(e^{\Xi_{t,\rho}}-1\right)^{1/\rho}.}
\tag{27}
\]

For score-gap pruning, replace \(r\) in the coupling lift by
\(e^{-\Delta}r\).  Lemma 3 contributes \(e^{-\Delta/2}\), and the exact
gap theorem from the companion note gives
\[
\boxed{
 D_t^{\rm gap}
 \le
 e^{-\Delta/2}\left(e^{\Xi_{t,1}}-1\right).}
\tag{28}
\]
The excess logical-Bayes-risk result becomes
\[
\boxed{
 R_{\rm Frontier}-R_{\rm logical\ ML}
 \le
 \sum_t\left[
 e^{-\Delta/2}(e^{\Xi_{t,1}}-1)
 +
 K^{-1/\rho}(e^{\Xi_{t,\rho}}-1)^{1/\rho}
 \right].}
\tag{29}
\]

\section{Connected-set and lifetime bounds}

Let \(d_t\) be the maximum number of processed variables incident to a
completed detector, with \(d_t=1\) by convention when there are no completed
incidences.  In the graphlike case the maximum degree of \(\mathcal G_t\) is
at most
\[
 \Delta_t=2(d_t-1).
\tag{30}
\]
A graph of maximum degree \(\Delta_t\) has at most
\(\Delta_t^{2(m-1)}\) connected \(m\)-sets containing a fixed root: encode a
canonical rooted spanning tree by its depth-first walk of length
\(2(m-1)\).  Therefore, with \(n_t=|I_t|\) and
\(\overline\beta_t=\max_{j\in I_t}\beta_j\),
\[
 \Xi_{t,\rho}
 \le
 n_t\sum_{m\ge1}
 \Delta_t^{2(m-1)}\overline\beta_t^{\rho m}.
\tag{31}
\]
If \(\Delta_t^2\overline\beta_t^\rho<1\), this gives
\[
\boxed{
 \Xi_{t,\rho}
 \le
 \frac{n_t\overline\beta_t^\rho}
 {1-\Delta_t^2\overline\beta_t^\rho}.}
\tag{32}
\]
For \(\Delta_t=0\), the \(m=1\) factor in (31) is interpreted as one.
This deliberately coarse count drops completed-parity and visibility
restrictions.

Finally, define
\[
 \mathcal L_\rho=\sum_t\Xi_{t,\rho},
 \qquad
 x_\rho=\sup_t\Xi_{t,\rho}.
\tag{33}
\]
For \(0<x_\rho<\infty\),
\[
 (e^x-1)^{1/\rho}
 \le
 e^{x_\rho/\rho}x_\rho^{1/\rho-1}x,
 \qquad 0\le x\le x_\rho.
\]
Equation (27) yields the integrated ordering-lifetime form
\[
\boxed{
 D^{\rm cap}
 \le
 K^{-1/\rho}
 e^{x_\rho/\rho}
 x_\rho^{1/\rho-1}
 \mathcal L_\rho.}
\tag{34}
\]
This isolates the time-integrated visible-polymer activity rather than the
worst-case raw boundary width.

\section{Scope and next steps}

The graphlike quotient-overlap family theorem is now proved, including
arbitrary quotient merging, invisible components, row-score effects, and a
fully explicit connected-set bound.  The exact low-weight overlap advantage
survives in the nonmultiplicative family bound (15), whereas the unconditional
componentwise theorem uses
\(\beta(p)=2\sqrt{p(1-p)}\sim2\sqrt p\).

The retained surface and BB144/Gross DEMs have detector-column weights four
and six, so this graphlike theorem does not apply to them directly.  The
bounded-hypergraph companion closes that applicability gap, computes exact
size-one/two ordering lifetimes, and shows that the resulting product-Peierls
certificate is already vacuous at practical caps.  The next improvement must
preserve boundary-shift aggregation or quotient cancellation, rather than
merely count larger polymers.  Terminal logical ranking through
charged/projective sector distortion remains separate.  None of the present
bounds equates raw frontier width with logical error.

\section*{References}

\begin{itemize}
\item H. Finner, ``A Generalization of H\"older's Inequality and Some
Probability Inequalities,'' \emph{The Annals of Probability} 20(4),
1893--1901 (1992).
\url{https://doi.org/10.1214/aop/1176989534}
\item O. Fawzi, A. Grospellier, and A. Leverrier, ``Efficient decoding of
random errors for quantum expander codes,'' arXiv:1711.08351 (2017).
\url{https://arxiv.org/abs/1711.08351}
\end{itemize}

\end{document}
